\section{Conic Sections: From Geometric Conditions to Equations}

The previous chapter studied sets determined by first-degree equations. A line
could be constructed from a point and a direction, or characterised by a
perpendicularity condition. We now ask a different kind of question:

\begin{center}
\emph{What curves arise when distances between moving and fixed geometric
objects are constrained?}
\end{center}

Distance in Cartesian coordinates contains squares. When a geometric distance
condition is written algebraically and square roots are removed, second-degree
terms appear naturally. Quadratic equations are therefore not introduced as an
arbitrary next list of formulas; they arise from simple Euclidean questions.

\subsection{A circle: one fixed point and one fixed distance}

Fix a centre $C=(h,k)$ and a number $r>0$. A point $P=(x,y)$ belongs to the
circle precisely when
\[
d(P,C)=r.
\]
Using the Euclidean distance formula,
\[
\sqrt{(x-h)^2+(y-k)^2}=r.
\]
Both sides are nonnegative, so squaring preserves the condition:
\[
\boxed{(x-h)^2+(y-k)^2=r^2.}
\]
The familiar quadratic equation has therefore been derived from the geometric
definition of a circle.

\subsection{A parabola: equal distance from a point and a line}

Fix a point $F$, called the focus, and a line $D$, called the directrix. A
parabola is the set of points whose distance from $F$ equals their distance from
$D$.

Choose coordinates so that
\[
F=(0,p),
\qquad
D:\ y=-p,
\qquad p>0.
\]
For $P=(x,y)$, the two distances are
\[
d(P,F)=\sqrt{x^2+(y-p)^2}
\]
and
\[
d(P,D)=|y+p|.
\]
On the resulting parabola $y\geq-p$, so $|y+p|=y+p$. The defining equality gives
\[
\sqrt{x^2+(y-p)^2}=y+p.
\]
Squaring and cancelling equal terms,
\[
x^2+y^2-2py+p^2=y^2+2py+p^2,
\]
so
\[
\boxed{x^2=4py.}
\]
The equation is not the definition. It is the Cartesian form of the equal-distance
condition.

\begin{center}
\begin{tikzpicture}[scale=0.85,>=stealth]
    \draw[axis] (-4.0,0) -- (4.2,0) node[right] {$x$};
    \draw[axis] (0,-2.0) -- (0,5.0) node[above] {$y$};
    \draw[gray,thick] (-3.7,-1.2) -- (3.7,-1.2)
        node[right] {directrix $D$};
    \coordinate (F) at (0,1.2);
    \coordinate (P) at (2.4,1.2);
    \coordinate (H) at (2.4,-1.2);
    \draw[subset,domain=-3.0:3.0,samples=100] plot (\x,{0.2083*\x*\x});
    \draw[blue!70!black,line width=1pt] (F)--(P)
        node[midway,above] {$d(P,F)$};
    \draw[orange!85!black,line width=1pt] (P)--(H)
        node[midway,right] {$d(P,D)$};
    \fill[geometricSubsetColor] (F) circle (2pt) node[above left] {$F$};
    \fill[geometricSubsetColor] (P) circle (2pt) node[above right] {$P$};
\end{tikzpicture}
\end{center}

\subsection{Two foci: ellipse and hyperbola}

Fix two distinct points $F_1,F_2$.

\begin{definitionbox}
An ellipse is the set of points $P$ for which
\[
\boxed{d(P,F_1)+d(P,F_2)=2a}
\]
for a fixed constant $2a>d(F_1,F_2)$.
\end{definitionbox}

The sum condition keeps the moving point in a bounded region. It is the
principle behind the classical string construction of an ellipse.

\begin{definitionbox}
A hyperbola is the set of points $P$ for which
\[
\boxed{|d(P,F_1)-d(P,F_2)|=2a}
\]
for a fixed positive constant smaller than $d(F_1,F_2)$.
\end{definitionbox}

The difference condition produces two separated branches. After choosing axes
adapted to the symmetry of the two foci and eliminating square roots, these
conditions yield the standard equations developed below.

\subsection{Why these curves form one family}

A circle, ellipse, parabola, and hyperbola can also be obtained by intersecting
a double cone with planes of different inclination. This geometric origin gives
the name \emph{conic sections}. The distance definitions explain the curves in
the plane; the cone picture explains why the same four forms belong to one
family.

\begin{remarkbox}
The rest of the chapter studies how these geometric sets appear in Cartesian
equations. Standard forms, translations, rotations, and discriminants are tools
for recognising the object already present in the equation. They do not replace
the geometric definition of the curve.
\end{remarkbox}
